\section{Scope, Chronology, and Evidence}

\subsection{Identity, names, and range}

Subject: Pope Saint Gregory I, known as Gregory the Great (\latin{Gregorius
Magnus}; in Byzantine tradition, Gregory the Dialogist, \latin{Dialogos});
Roman urban prefect, Benedictine-era monastic founder, regional deacon, papal
\latin{apocrisiarius} at Constantinople, bishop of Rome (590--604), Church
Father, and one of the four original Latin Doctors of the Church. Whole-life
range: born at Rome around 540 (exact year unrecorded); died at Rome on 12
March 604; consecrated bishop of Rome on 3 September 590. Ordinary English
names are used throughout, with Latin technical and ecclesiastical terms
translated where they first appear.

\subsection{Corpus, citation conventions, and reporting ceilings}

Quotations of the \work{Regula Pastoralis} and \work{Selected Epistles} (Books
I--VIII) follow James Barmby's unrevised 1895 translation in \work{Nicene and
Post-Nicene Fathers}, Series 2, vol.~12 (tracked CCEL text), which follows the
Benedictine/Maurist numbering in Migne \work{Patrologia Latina} 77 (differing
in places by one or two numbers from the later MGH Ewald--Hartmann edition).
Quotations of the \work{Moralia in Iob} follow the Oxford Library of the
Fathers English translation (1844--1850), inspected directly in the tracked
complete text. Latin quotations of the \work{Homiliae XL in Evangelia} and
\work{Homiliae in Hiezechihelem Prophetam} were verified directly in tracked
Migne PL 76 and Wikisource Latin texts, with English renderings supplied inline
where no public-domain tracked English version stood beside them. Gregory of
Tours's \work{Historia Francorum} X.1 was inspected in both O.\ M.\ Dalton's
1927 Oxford translation and W.\ Arndt's 1884 MGH Latin edition. Liturgical
witnesses were inspected directly in H.\ A.\ Wilson's 1915 \work{Hadrianum}
edition, Migne PL 78, Adrian Fortescue's verified passages (1922, pp.\ 172,
234--235), and the 1902 \work{Martyrologium Romanum}.

Several important witnesses are not held in full among the editions consulted
for this study and are cited strictly through their reporting witnesses: Gregory's \work{Dialogi} and
\work{Registrum} Books IX--XIV (reported via Barmby 1895, Fortescue 1922, and
Migne PL 76/78); Bede's \work{Historia Ecclesiastica Gentis Anglorum}, the
anonymous Whitby \work{Vita Gregorii}, Paul the Deacon's \work{Vita} and
\work{Historia Langobardorum}, and John the Deacon's \work{Vita S. Gregorii}
(reported via Barmby 1895, Wilson 1915, Fortescue 1922, and Migne PL 78); and
Pope Boniface VIII's 1298 decretal on the four Latin Doctors (reported via
standard historical reference). No original archival, manuscript, or
archaeological fieldwork was conducted. Current liturgical calendar notes
reflect the 1969 General Roman Calendar alongside the 1902 Roman Martyrology as
of November 2026.

\subsection{Evidence classes}

\begin{evidencekey}
\evidenceclass{A}{Subject-authored evidence: Gregory's \work{Moralia in Iob}
(including Epistle Dedicatory and XIV.56), \work{Regula Pastoralis},
\work{Registrum Epistolarum}, \work{Homiliae XL in Evangelia}, and
\work{Homiliae in Hiezechihelem}; each claim qualified by genre, occasion, and
transmission.}
\evidenceclass{N}{Contemporary external witness: Gregory of Tours,
\work{Historia Francorum} X.1 (eyewitness report of the deacon of Tours in
Rome, 590); Licinianus of Cartagena (\work{Ep.}\ II.54).}
\evidenceclass{E}{Early received tradition: eighth-century English and Lombard
testimony (Whitby \work{Vita}, Bede \work{HE}, Paul the Deacon), establishing
early memory and cult rather than eyewitness fact.}
\evidenceclass{L}{Later hagiographic, liturgical, iconographic, and medieval
tradition: John the Deacon (late 9th c.), Carolingian sacramentary and chant
attributions, the Castel Sant'Angelo vision, the whispering dove, formalization
of Gregorian Trentals, and the Trajan legend.}
\evidenceclass{M}{Official ecclesial reception, limited to each act's actual
object: \work{Hadrianum} feast of 12 March, 1902 \work{Martyrologium Romanum}
(12 March and 3 September), Boniface VIII's 1298 Doctors decretal (reported),
and the 1969 General Roman Calendar memorial on 3 September.}
\evidenceclass{K}{Named modern critical reconstruction: James Barmby (1895),
H.\ A.\ Wilson (1915), Adrian Fortescue (1922), and modern patristic and
musicological consensus on \work{Dialogi} authorship and Frankish-Roman chant.}
\evidenceclass{S}{Source-grounded historical synthesis, distinct from any
single witness.}
\end{evidencekey}

\subsection{Chronology}

\begin{lifetimeline}
\lifeevent{c.~540}{Birth at Rome into a wealthy senatorial Christian family;
descendant of Pope Felix III (483--492), whom Gregory calls his \latin{atavus}
(``forefather'')}{A (\work{Hom.\ in Evang.}\ 38.15); N (\work{HF} X.1); K
(Barmby)}{Year approximate; lineage secure}
\lifeevent{568}{Lombards under Alboin invade Italy, shaping the political world
of Gregory's adulthood}{K (Barmby Prolegomena, citing Latin \latin{viginti
septem} in \work{Ep.}\ V.21 and XIII.38; NPNF2-12 \work{Ep.}\ V.21 prints
``seven years'')}{Secure}
\lifeevent{by 573}{Serves as Prefect of the City of Rome (\latin{Praefectus
Urbi}) under Emperor Justin II}{N (\work{HF} X.1, silk attire); K (Barmby)}
{Secure}
\lifeevent{c.~574--575}{Monastic conversion after father's death; endows six
monasteries in Sicily and Saint Andrew's on the Caelian Hill near SS.\ John and
Paul}{A (\work{Moralia}, Ep.\ Ded.\ 1; \work{Hom.\ in Evang.}\ 38.16); N
(\work{HF} X.1)}{Episode secure; exact year approximate}
\lifeevent{c.~578--579}{Ordained seventh regional deacon of Rome (by Benedict I
or Pelagius II) and sent to Constantinople as papal \latin{apocrisiarius}}{A
(\work{Moralia}, Ep.\ Ded.\ 1); N (\work{HF} X.1); K (Barmby)}{Event secure;
ordaining pope debated}
\lifeevent{c.~579--585}{Resides in Constantinople; friendship with Leander of
Seville; begins oral conferences of \work{Moralia in Iob}; disputes with
Patriarch Eutychius (d.~582) before Emperor Tiberius II; godfather to Maurice's
son}{A (\work{Moralia}, Ep.\ Ded.\ 1; XIV.56); N (\work{HF} X.1)}{Secure}
\lifeevent{c.~585--586}{Recalled to Rome by Pelagius II; returns to Saint
Andrew's monastery}{K (Barmby); E (Paul the Deacon, reported)}{Approximate}
\lifeevent{Nov.\ 589--Feb.\ 590}{Tiber flood ruins granaries; \latin{lues
inguinaria} plague strikes Rome; Pope Pelagius II dies (7/8 Feb.\ 590); Gregory
elected and appeals to Maurice against confirmation}{N (\work{HF} X.1)}{Secure}
\lifeevent{Spring--Aug.\ 590}{Preaches \work{Oratio de mortalitate} and leads
the sevenfold litany (\latin{litania septiformis}) to Saint Mary Major}{N
(\work{HF} X.1); A (PL 76, col.\ 1311)}{Secure}
\lifeevent{3 Sept.\ 590}{Consecrated bishop of Rome in Saint Peter's Basilica
after attempted flight}{N (\work{HF} X.1); M (\work{Mart.\ Rom.}\ 1902)}
{Secure}
\lifeevent{590--591}{Composes \work{Regula Pastoralis}; issues synodical
confession (\work{Ep.}\ I.25); reforms Sicilian patrimony (\work{Ep.}\ I.44);
forbids forced baptism of Jews (\work{Ep.}\ I.47)}{A (\work{Reg.\ Past.};
\work{Ep.}\ I.5, 10, 25, 35, 43--44, 47)}{Secure}
\lifeevent{590--592}{Preaches \work{Homiliae XL in Evangelia} in Rome's
stational basilicas; defends Rome against Duke Ariulf of Spoleto}{A
(\work{Hom.\ in Evang.}; \work{Ep.}\ II.46)}{Secure}
\lifeevent{593}{Preaches \work{Homiliae in Hiezechihelem}; halts exposition
during King Agilulf's siege of Rome; protests Maurice's monastic law
(\work{Ep.}\ III.65)}{A (\work{Hom.\ in Hiez.}\ II.6, 10; \work{Ep.}\ III.65,
V.40)}{Secure}
\lifeevent{593--594}{Composes \work{Dialogi} (including Book II on Saint
Benedict and Book IV on the afterlife)}{A/K (reported via Barmby and PL 76)}
{Secure as span}
\lifeevent{595}{Sends completed \work{Moralia} and \work{Regula Pastoralis} to
Leander (\work{Ep.}\ V.49); launches campaign against ``Universal Bishop''
title (\work{Ep.}\ V.18--21, 43; dated Jan.\ 595 in Barmby's Benedictine text
of V.18, June 595 in MGH); instructs Candidus to buy English youths in Gaul
(\work{Ep.}\ VI.7)}{A (\work{Ep.}\ V.18--21, 40--43, 49; VI.7)}{Secure}
\lifeevent{596--598}{Dispatches Augustine of Canterbury to England (23 July
596, \work{Ep.}\ VI.51--59); announces 10,000 English baptisms at Christmas 597
and refuses ``Universal Pope'' (\work{Ep.}\ VIII.30)}{A (\work{Ep.}\ VI.51--59,
VIII.30)}{Secure}
\lifeevent{598--599}{Lombard truce concluded via Queen Theodelinda; defends
Roman liturgical customs in \work{Ep.}\ IX.12}{A/K (Fortescue 1922; Barmby)}
{Secure}
\lifeevent{601}{Sends Mellitus, Justus, Paulinus, pallium, \work{Libellus
Responsionum}, and temple-adaptation policy to England}{A/E (reported via
Barmby from \work{Ep.}\ XI.28--29, 64--66, 76 and Bede \work{HE} I.27--32)}
{Secure}
\lifeevent{602--603}{Usurper Phocas overthrows Maurice (Nov.\ 602); Gregory
writes congratulatory letters (603); renewed Lombard truce (Sept.\ 603)}{A/K
(reported via Barmby from \work{Ep.}\ XIII.31, 38--39; XIV.12)}{Secure}
\lifeevent{12 Mar.\ 604}{Death at Rome; burial in the portico of Old Saint
Peter's Basilica}{M (\work{Hadrianum}; \work{Mart.\ Rom.}\ 1902); K (Barmby)}
{Secure}
\lifeevent{c.~704--731}{Whitby \work{Vita} and Bede's \work{Historia
Ecclesiastica} establish early English cult and record \latin{Non Angli} and
Trajan legends}{E (reported via Barmby)}{Secure as textual dates}
\lifeevent{784--791}{Pope Hadrian I sends the \work{Hadrianum} Sacramentary to
Charlemagne; Carolingian \latin{Hucusque} Supplement distinguishes Gregory's
reputed core from later additions}{L/K (Wilson 1915; Migne PL 78)}{Secure}
\lifeevent{c.~873--875}{John the Deacon writes four-book \work{Vita S.
Gregorii}, attributing the \work{Antiphonarius} and \latin{schola cantorum}
to Gregory}{L (quoted in Migne PL 78; reported via Barmby)}{Secure}
\lifeevent{1298}{Pope Boniface VIII confirms the feasts of the four original
Latin Doctors of the Church}{M (reported)}{Secure}
\end{lifetimeline}

\subsection{Principal works cited}

\begin{worksinventory}
\workentry{\work{Moralia in Iob} (35 books)}{c.~579--595}{Begun orally in
Constantinople; dedicated to Leander of Seville}{Constantinople / Pontificate}
{Ep.\ Ded.\ 1--5; Pref.\ 2; XIV.56}
\workentry{\work{Regula Pastoralis} (4 parts)}{590--591}{Addressed to John of
Ravenna on the weight of pastoral office}{Early Pontificate}{I.1--7; II.1,
4--6; III Prol., 1; IV}
\workentry{\work{Oratio de mortalitate}}{590}{Plague sermon ordering the
sevenfold litany in Rome}{Pope-elect}{\work{HF} X.1; PL 76, col.\ 1311}
\workentry{\work{Homiliae XL in Evangelia}}{590--592}{Stational preaching in
Rome; dedicated to Secundinus of Taormina}{Early Pontificate}{Praef.; 1.1, 1.5;
8.1; 20; 37.7--9; 38.15--16; 40.7, 11--12}
\workentry{\work{Homiliae in Hiezechihelem} (22 hom.)}{593}{Preached in Rome;
interrupted by Agilulf's siege}{Middle Pontificate}{II.6.1, 22; II.7.7--9;
II.10.24}
\workentry{\work{Registrum Epistolarum} (14 books)}{590--604}{Papal governance,
patrimonies, diplomacy, missions, spiritual counsel}{Pontificate}{I--VIII
(direct); IX--XIV (reported)}
\workentry{\work{Dialogi} (4 books)}{593--594}{Italian holy men, Saint
Benedict, afterlife, and Mass for the dead}{Middle Pontificate}{II; III.33;
IV.55/57 (reported)}
\end{worksinventory}

\subsection{Summary of audited traditions}

\begin{traditionaudit}
\traditionrow{\latin{Non Angli sed Angeli} slave-market puns}{Whitby
\work{Vita} (early 8th c.); Bede \work{HE} II.1 (c.~731)}{Absent from Gregory's
letters; \work{Ep.}\ VI.7 shows Gallic slave purchases in 595}{Early received
English tradition (E), distinct from contemporary evidence}
\traditionrow{Archangel Michael sheathing sword atop Hadrian's Mausoleum
(590)}{Medieval Roman legend / \work{Golden Legend} (13th c.)}{Absent from
eyewitness Gregory of Tours \work{HF} X.1 and from Bede}{Later devotional
legend (L); unattested in 590}
\traditionrow{White dove whispering at Gregory's ear during dictation}{Paul the
Deacon / John the Deacon (8th--9th c.)}{Gregory's own prefaces describe
ordinary dictation, illness, and scribal revision}{Later hagiographic and
iconographic emblem (L)}
\traditionrow{``Gregorian Mass'' (thirty consecutive daily Masses)}
{\work{Dialogi} IV.55/57 (593--594, case of monk Justus)}{One monastic case in
Gregory; formalized as universal custom only in Middle Ages}{Authentic
Gregorian origin (A, reported); medieval formalization (L)}
\traditionrow{Intercession delivering pagan Emperor Trajan from hell}{Whitby
\work{Vita} ch.\ 29 (early 8th c.); Paul and John the Deacon}{Absent from
Gregory's corpus; in tension with \work{Hom.\ in Evang.}\ 40.7}{Early medieval
legend (E/L) used in scholastic theology}
\traditionrow{Personal authorship of complete Gregorian Sacramentary and chant
corpus}{Carolingian rubrics (8th--9th c.); John the Deacon \work{Vita} II.6
(c.~873--875)}{\work{Hadrianum} includes Gregory's own feast and \latin{Hucusque}
Supplement; chant crystallized in 8th--9th c.}{Later attribution (L) overlying
genuine specific reforms (A/N)}
\end{traditionaudit}

\subsection{Rights and review state}

Original narrative prose is project-authored. Quotations of Gregory's
\work{Regula Pastoralis} and \work{Epistles} (Barmby 1895, NPNF2-12), the
\work{Moralia in Iob} (Library of the Fathers, 1844--1850), Gregory of Tours
(Dalton 1927; Arndt 1884), Migne PL 76 and PL 78, Wilson's \work{Hadrianum}
(1915), Fortescue (1922), and the 1902 \work{Martyrologium Romanum} are drawn
from public-domain texts retained in the source library. Inline English
translations of tracked Latin passages from the \work{Homiliae in Evangelia}
and \work{Homiliae in Hiezechihelem} are project renderings of public-domain
Latin. Internal source, editorial, web-edition, and print checks have been
completed; no independent specialist, theological, canonical, or ecclesiastical
approval is claimed.
